\documentclass[a4paper,12pt]{article}
\usepackage[utf8]{inputenc}
\usepackage[T1]{fontenc}
\usepackage[norsk]{babel}
\usepackage{amsmath}
\usepackage{amsthm}
\usepackage{listings}
\usepackage{xcolor}
\usepackage{geometry}
%\geometry{margin=2.5cm}
\usepackage{enumitem}
\usepackage{lmodern}
\usepackage{cmap}
\usepackage{needspace}
\usepackage[colorlinks=true, linkcolor=red!50!black, citecolor=red!50!black, urlcolor=blue]{hyperref}
\usepackage[all]{hypcap}
\usepackage[stable]{footmisc}

\usepackage{setspace}
\onehalfspacing

\usepackage{tikz}
\usepackage{tikz-3dplot}
\usetikzlibrary{angles,quotes}

\renewcommand{\qedsymbol}{QED}


\title{Areal og volum med vektorregning}
\date{}

\begin{document}

\maketitle

\section*{Oppgave 1: Parallellogram i 2D}
Vi har et parallellogram utspent av to vektorer $\vec{u}$ og $\vec{v}$. Vinkelen mellom vektorene er $\alpha$.

\begin{figure}[h]
\begin{center}
\begin{tikzpicture}[scale=1.2]
    \draw[->, thick] (0,0) -- (3,0) node[midway, below] {$\vec{u}$};
    \draw[->, thick] (0,0) -- (1,2) node[midway, left] {$\vec{v}$};
    \draw[dashed] (3,0) -- (4,2) -- (1,2);
    \draw[red, thick] (1,0) -- (1,2) node[midway, right] {$h$};
    \draw (0.4,0) arc (0:63.4:0.4) node[pos=0.5, right] {$\alpha$};
    \node at (2.5,1.5) [below left] {$A$};
\end{tikzpicture}
\caption{Parallellogram utspent av to vektorer.}
\end{center}
\end{figure}

Vis at arealet $A$ av parallellogrammet er gitt ved 
\begin{equation}
	A = |\vec{u}| |\vec{v}| \sin(\alpha).
\end{equation}


\section*{\boldmath Oppgave 2: Areal og $2\times2$-determinant}

I 2 dimensjoner er determinanten av to vektorer $\vec{u} = [u_1, u_2]$ og $\vec{v} = [v_1, v_2]$ gitt ved \[ \det(\vec{u}, \vec{v}) = \begin{vmatrix} u_1 & u_2 \\ v_1 & v_2 \end{vmatrix} = u_1v_2 - u_2v_1. \]

\begin{enumerate}
    \item[a)] Gitt $\vec{u} = [4, 1]$ og $\vec{v} = [1, 3]$, finn determinanten $\det(\vec u, \vec v)$.
    \item[b)] Sammenlign verdien av determinanten med arealet av parallellogrammet vektorene utspenner. Dette kan du gjøre i geogebra eller for hånd.
\end{enumerate}

Poenget er at i 2 dimensjoner så gir determinanten samme svar som arealet. Dette er alltid tilfellet og vi kan derfor si at
\begin{equation}
\label{eq:2dlikhet}
|\vec{u}| |\vec{v}| \sin(\alpha) =  u_1v_2 - u_2v_1.
\end{equation}

\section*{Oppgave 3: Kryssprodukt og areal i 3D}
I 3 dimensjoner definerer vi kryssproduktet $\vec{u} \times \vec{v}$ ved hjelp av en determinant:
\[
\vec{u} \times \vec{v} = \begin{vmatrix} \vec{e}_x & \vec{e}_y & \vec{e}_z \\ u_1 & u_2 & u_3 \\ v_1 & v_2 & v_3 \end{vmatrix}
\]
En viktig egenskap er at $|\vec{u} \times \vec{v}| = |\vec{u}| |\vec{v}| \sin(\alpha)$.

\begin{enumerate}
    \item[a)] La $\vec{u} = [4, 1, 0]$ og $\vec{v} = [1, 3, 0]$. Regn ut $\vec{u} \times \vec{v}$.
    \item[b)] Hva blir lengden av denne vektoren, og hvordan samsvarer dette med det du lærte i Oppgave 1 og 2?
	\item[c)] Finn plass i hodet ditt til å huske at $|\vec u \times \vec v|$ gir arealet til parallellogrammet utspent av $\vec u$ og $\vec v$.
	\item[d)] Hva blir arealet av trekanten utspent av $\vec u$ og $\vec v$?
\end{enumerate}

\newpage

\section*{\boldmath Oppgave 4: Volum og $3\times 3$-determinant}

Vi har sett at determinanten i 2D gir arealet av parallellogrammet spent ut av to vektorer. I 3D gir determinanten oss volumet $V$ av \textit{parallellepipedet} (figur \ref{fig:parallellepiped}) mellom tre vektorer $\vec{u}, \vec{v}$ og $\vec{w}$:\phantomsection\footnote{I 4D gir determinanten oss \textit{hyper}-volumet mellom fire vektorer osv.}

\[
V = \det(\vec w,\vec u,\vec v) = \begin{vmatrix} w_1 & w_2 & w_3 \\ u_1 & u_2 & u_3 \\ v_1 & v_2 & v_3 \end{vmatrix}.
\]

\begin{figure}[h]
\centering
\begin{minipage}{0.48\textwidth}
\tdplotsetmaincoords{70}{110}
\begin{tikzpicture}[tdplot_main_coords, scale=1.5]

% --- Define vectors ---
\coordinate (u) at (2,0,0);
\coordinate (v) at (0,2,0);
\coordinate (w) at (1.1,1.1,2);

% --- Main vectors ---
\draw[->, thick] (0,0,0) -- (u) node[anchor=north east]{$\vec{u}$};
\draw[->, thick] (0,0,0) -- (v) node[anchor=north west]{$\vec{v}$};
\draw[->, thick] (0,0,0) -- (w) node[anchor=south]{$\vec{w}$};
\draw[blue,->, thick] (0,0,0) -- (0,0,2.5) node[anchor=east]{$\vec{u} \times \vec v $};

% --- Helper lines ---
\draw[red, thick] (w) -- ($(w)+(0,0,-2)$) node[midway, right] {$h$};
\draw[dashed] (0,0,0) -- ($(w)+(0,0,-2)$);

\coordinate (A) at ($0.6*0.44*(w)$);
\coordinate (B) at ($0.6*(0.7,0.7,0)$);

\draw (A) to[bend left=25] node[midway,right] {$\theta$} (B);

\draw[dashed] (u) -- ($(u)+(v)$) -- (v);
\draw[dashed] (w) -- ($(w)+(u)$) -- ($(w)+(u)+(v)$) -- ($(w)+(v)$) -- cycle;

\draw[dashed] (u) -- ($(u)+(w)$);
\draw[dashed] ($(u)+(v)$) -- ($(u)+(v)+(w)$);
\draw[dashed] (v) -- ($(v)+(w)$);

\end{tikzpicture}
\end{minipage}
\hfill
\begin{minipage}{0.48\textwidth}
\centering
\begin{tikzpicture}

% Example 2D diagram
\coordinate (O) at (0,0);
\coordinate (X) at (1.6,0);
\coordinate (W) at (1.6,2);

% Example 2D diagram
\coordinate (O) at (0,0);
\coordinate (X) at (1.6,0);
\coordinate (W) at (1.6,2);

% Example 2D diagram
\draw[dashed] (O) -- (X);
\draw[blue,->, thick] (O) -- (0,4) node[anchor=east] {$\vec u \times \vec v$};
\draw[->, thick] (O) -- (W) node[anchor=south] {$\vec w$};
\draw[red, thick] (W) -- (X) node[midway, right] {$h$};

% Angle arc
\draw pic[draw, angle radius=0.8cm] {angle = X--O--W};

% Label
\node at (0.5,0.25) {$\theta$};

\end{tikzpicture}
\end{minipage}
\caption{Parallellepiped utspent av $\vec{u}, \vec{v}$ og $\vec{w}$. Vinkelen mellom $\vec w$ og $uv$-planet er $\theta$.}
\label{fig:parallellepiped}
\end{figure}


\begin{enumerate}[label=\alph*)]
	\item En annen måte å skrive denne determinanten på er $(\vec u \times\vec v) \cdot \vec w$. Vis at denne gir rett formel for volumet av parallellepipedet:\footnote{Hint: Finn uttrykk for både $h$ og $(\vec u \times\vec v) \cdot \vec w$ ved hjelp av $\theta$.} \[V=Ah,\]der $A$ er arealet av grunnflaten og $h$ er høyden.
    \item Regn ut volumet av parallellepipedet utspent av $\vec{u}=[3,0,0]$, $\vec{v}=[0,4,0]$ og $\vec{w}=[1,1,5]$.
	\item Regn ut volumet av parallellepipedet utspent av $\vec{u}=[3,0,0]$, $\vec{v}=[0,4,0]$ og $\vec{w}=[-1,-1,-5]$. Hvorfor tror du svaret nå ble negativt?\footnote{Bruk høyerehåndsregelen.}
\end{enumerate}

\section*{Bonus: Vis at determinanten gir areal i 2D}

Du skal altså vise dette:
\[\det(\vec u, \vec v)= u_1v_2 - u_2v_1 = |\vec{u}| |\vec{v}| \sin(\alpha).\]

\begin{enumerate}

    \item[a)] Begynn med å kvadrere determinanten og vis at den kan skrives slik:
	\[
		(\det(\vec u, \vec v))^2 = (u_1^2 +u_2^2)(v_1^2 +v_2^2) - (u_1v_1 + u_2v_2)^2.
	\]
    \item[b)] Sett inn uttrykk for skalarproduktet på koordinatform: $|\vec u|^2$, $|\vec v|^2$ og $\vec u \cdot \vec v$. Vis at vi nå kan si:
	\[
		(\det(\vec u, \vec v))^2 = |\vec u|^2|\vec v|^2 - (\vec u \cdot \vec v)^2.
	\]
	    \item[c)] Bruk uttrykket for skalarproduktet med vinkel og vis at:
	\[
		(\det(\vec u, \vec v))^2 = |\vec u|^2|\vec v|^2(1-\cos^2(\alpha)).
	\]
	    \item[d)] Fullfør beviset ved å vise at:
	\[
		\det(\vec u, \vec v) = |\vec u||\vec v|\sin(\alpha).
		\tag*{\qed}
	\]

\end{enumerate}


\appendix
\newpage
\section{Hvorfor gir kryssproduktet areal?}

Vi vet fra geometri at arealet av et parallellogram spent ut av to vektorer er

\[
A = |\vec u|\,|\vec v|\sin\alpha
\]

der $\alpha$ er vinkelen mellom vektorene.

Vi skal nå vise at dette er det samme som lengden av kryssproduktet.

\subsection*{Steg 1: Kryssprodukt i koordinatform}

La

\[
\vec u =
\begin{pmatrix}
u_1\\
u_2\\
u_3
\end{pmatrix},
\qquad
\vec v =
\begin{pmatrix}
v_1\\
v_2\\
v_3
\end{pmatrix}.
\]

Kryssproduktet er

\[
\vec u \times \vec v =
\begin{pmatrix}
u_2v_3-u_3v_2 \\
u_3v_1-u_1v_3 \\
u_1v_2-u_2v_1
\end{pmatrix}.
\]

Lengden av denne vektoren er

\[
|\vec u \times \vec v|
=
\sqrt{
(u_2v_3-u_3v_2)^2
+
(u_3v_1-u_1v_3)^2
+
(u_1v_2-u_2v_1)^2.
}
\]

\subsection*{Steg 2: Kvadrer lengden}

Vi ser først på kvadratet av lengden:

\[
|\vec u \times \vec v|^2
=
(u_2v_3-u_3v_2)^2
+
(u_3v_1-u_1v_3)^2
+
(u_1v_2-u_2v_1)^2.
\]

Hvis vi utvikler kvadratene og samler leddene, får vi

\[
|\vec u \times \vec v|^2
=
(u_1^2+u_2^2+u_3^2)(v_1^2+v_2^2+v_3^2)
-
(u_1v_1+u_2v_2+u_3v_3)^2.
\]

\subsection*{Steg 3: Bruk prikkproduktet}

Vi kjenner to viktige størrelser:

Lengden av vektorene

\[
|\vec u|^2 = u_1^2+u_2^2+u_3^2
\]

\[
|\vec v|^2 = v_1^2+v_2^2+v_3^2,
\]

og prikkproduktet

\[
\vec u\cdot\vec v
=
u_1v_1+u_2v_2+u_3v_3
=
|\vec u||\vec v|\cos\alpha.
\]

Dermed kan vi skrive

\[
|\vec u \times \vec v|^2
=
|\vec u|^2 |\vec v|^2
-
(\vec u\cdot\vec v)^2.
\]

Setter vi inn uttrykket for prikkproduktet får vi
\begin{align*}
|\vec u \times \vec v|^2 
&= |\vec u|^2 |\vec v|^2 - (|\vec u||\vec v|\cos\alpha)^2 \\
&= |\vec u|^2 |\vec v|^2 (1-\cos^2\alpha).
\end{align*}

\subsection*{Steg 4: Bruk trigonometrisk identitet}

Vi bruker identiteten

\[
1-\cos^2\alpha = \sin^2\alpha.
\]

Dermed får vi

\[
|\vec u \times \vec v|^2
=
|\vec u|^2 |\vec v|^2 \sin^2\alpha.
\]

Tar vi kvadratroten får vi

\[
|\vec u \times \vec v|
=
|\vec u|\,|\vec v|\,\sin\alpha. \tag*{\qed}
\]

\subsection*{Konklusjon}

Lengden av kryssproduktet er altså nøyaktig lik arealet av parallellogrammet spent ut av vektorene:

\[
A = |\vec u \times \vec v|.
\]

Dette forklarer hvorfor kryssproduktet brukes til å finne areal i rommet.
\newpage
\section{Hvorfor gir skalarproduktet vinkelen mellom to vektorer?}

Skalarproduktet kan beregnes på to måter:

\[
\vec u \cdot \vec v = u_1v_1 + u_2v_2 + u_3v_3
\]

og

\[
\vec u \cdot \vec v = |\vec u|\,|\vec v|\cos\alpha,
\]

der $\alpha$ er vinkelen mellom vektorene. Vi skal vise hvorfor disse to uttrykkene er like.

\subsection*{Steg 1: Start med lengden av $\vec u-\vec v$}

La

\[
\vec u =
\begin{pmatrix}
u_1\\
u_2\\
u_3
\end{pmatrix},
\qquad
\vec v =
\begin{pmatrix}
v_1\\
v_2\\
v_3
\end{pmatrix}.
\]

Da er

\[
\vec u - \vec v =
\begin{pmatrix}
u_1-v_1\\
u_2-v_2\\
u_3-v_3
\end{pmatrix}.
\]

Lengden av denne vektoren er

\[
|\vec u-\vec v|^2 =
(u_1-v_1)^2 + (u_2-v_2)^2 + (u_3-v_3)^2.
\]

\subsection*{Steg 2: Utvid uttrykket}

Utvider vi kvadratene får vi

\[
|\vec u-\vec v|^2 =
u_1^2+u_2^2+u_3^2
+
v_1^2+v_2^2+v_3^2
-
2(u_1v_1+u_2v_2+u_3v_3).
\]

Vi kjenner igjen

\[
|\vec u|^2 = u_1^2+u_2^2+u_3^2
\]

\[
|\vec v|^2 = v_1^2+v_2^2+v_3^2.
\]

Dermed får vi

\[
|\vec u-\vec v|^2 =
|\vec u|^2 + |\vec v|^2 - 2(u_1v_1+u_2v_2+u_3v_3).
\]

\subsection*{Steg 3: Sammenlign med cosinussetningen}

Hvis vi ser på trekanten som dannes av $\vec u$, $\vec v$ og $\vec u-\vec v$, gir cosinussetningen

\[
|\vec u-\vec v|^2 =
|\vec u|^2 + |\vec v|^2 - 2|\vec u||\vec v|\cos\alpha.
\]

\begin{figure}[h]
\centering
\begin{tikzpicture}[scale=1.5]
% Vectors
\coordinate (O) at (0,0);
\coordinate (U) at (2,0.5);
\coordinate (V) at (1,1.8);

% Draw vectors u and v
\draw[->, thick, blue] (O) -- (U) node[midway, below] {$\vec u$};
\draw[->, thick, red] (O) -- (V) node[midway, left] {$\vec v$};

% Draw u - v
\draw[->, thick, green!60!black] (V) -- (U) node[midway, above right] {$\vec u-\vec v$};

% Angle alpha
\pic[draw, angle radius=0.6cm] {angle = U--O--V};
\node at (0.45,0.3) {$\alpha$};
\end{tikzpicture}

\end{figure}

\subsection*{Steg 4: Sammenlign uttrykkene}

Vi har nå to uttrykk for $|\vec u-\vec v|^2$:

\[
|\vec u|^2 + |\vec v|^2 - 2(u_1v_1+u_2v_2+u_3v_3)
\]

og

\[
|\vec u|^2 + |\vec v|^2 - 2|\vec u||\vec v|\cos\alpha.
\]

Dermed må

\[
u_1v_1+u_2v_2+u_3v_3
=
|\vec u||\vec v|\cos\alpha. \tag*{\qed}
\]

\subsection*{Konklusjon}

Skalarproduktet kan derfor skrives på to måter:

\[
\vec u\cdot\vec v
=
u_1v_1+u_2v_2+u_3v_3\,\,\,\,
\land\,\,\,\,
\vec u\cdot\vec v
=
|\vec u||\vec v|\cos\alpha.
\]

Dette gjør det mulig å finne vinkelen mellom to vektorer:

\[
\cos\alpha =
\frac{\vec u\cdot\vec v}{|\vec u||\vec v|}.
\]

\end{document}